% ---------------------------------------------------------------------------
% Response to R1 Major 4 (all four parts) and R2 Major 4 -- JEDC revision
% Body only: paste after the previous comment in the combined response file.
% Needs booktabs in the preamble. Headings are \section*/\subsection*; if the
% other comments use \section* internally, promote these one level.
% ---------------------------------------------------------------------------

\section*{R1, Major 4 \textnormal{\itshape (and R2, Major 4)}}
{\itshape What is specifically ``oil'' about the mechanism?}

\begin{quote}\itshape
Mathematically, the same structure appears to apply to volatility in the price of
any sufficiently important imported intermediate input. [\ldots] the paper would
benefit from sharpening the discussion of what the oil-specific production
structure adds relative to these existing mechanisms.
\end{quote}

The referee's reading of the structure is correct and we now say so in the text:
mathematically this is an imported-input mechanism, not an oil mechanism. Nothing
in the CES technology knows that the input is oil. What follows sets out how we now
frame the contribution, and what the model and the data say about how much of the
effect is specific to oil.

\subsection*{Sharpening the contribution}

We take this point and have rewritten the discussion around it. The gap
statement in the submitted version was framed more broadly than the argument
requires, and the revision states it precisely: what is distinctive here is the
combination of oil \emph{importers}, an oil price process whose persistence and
stochastic volatility are \emph{estimated} rather than assumed, and a
transmission that runs through the \emph{cost} side of production. Section 2 now
discusses Lopez-Martin et al.\ (2019) rather than citing it in passing, and
brings the exporter studies of Hamann et al.\ (2018) and Esquivel (2020) from
footnote 1 into the text, where the contrast with an importing economy is the
point.

We are also clearer about what the paper builds on rather than introduces. The
stochastic-volatility machinery is standard, and long-duration debt follows
Chatterjee and Eyigungor (2012) and Hatchondo and Martinez (2009). State
dependence in volatility and debt is established by Johri, Khan and
Sosa-Padilla (2022) for world interest rate volatility; our Figures 3 and 5
display the same qualitative pattern for a different shock, and the revision
presents them as corroborating that finding in a new setting.

\subsection*{Differentiation from the closest papers}

\textbf{Lopez-Martin et al.\ (2019)}, and the exporter studies of Hamann et al.\
(2018) and Esquivel (2020), treat the commodity as a source of \emph{revenue};
here it is a \emph{cost} of production. That is a difference in the sign of the
budget constraint rather than a relabelling: a price rise is a windfall there and
a squeeze here, so default risk moves in the opposite direction.

\textbf{Seoane (2019)} studies stochastic volatility in the endowment. The
volatile object is exogenous income, so there is no optimizing margin. In our
model the oil price enters through the CES first-order condition, so a price
change moves the imported quantity, the import bill and output together, and
because our estimated elasticity is below one the oil bill \emph{rises} as a
share of GDP when oil is dear, with $d\ln s/d\ln p = +0.61$ (s.e.\ 0.08) in our
import panel. More importantly, the channel differs in sign. Because
$\mathrm{E}[p'] = \exp(m + v/2)$, higher volatility raises the expected
\emph{level} of the oil price, and our decomposition puts 12.1 of the 13.4 basis
points attributable to time-varying volatility in that expected-price channel,
with only 1.3 in pure dispersion. The identical convexity in an
endowment-volatility model raises expected \emph{income}.

\textbf{Johri, Khan and Sosa-Padilla (2022)} study volatility in the world
interest rate, which acts on the financing side and leaves production untouched.
Oil price volatility acts on production first and reaches the bond price only
through repayment capacity. The two are not nested.

\subsection*{How much of the quantitative effect is specific to oil}

\begin{quote}\itshape
If feasible, a simple model-based comparison with a generic imported-input price
process matched on persistence and unconditional volatility could also help
assess how much of the quantitative effect is specific to oil.
\end{quote}

This comparison is already in the paper, and we had not presented it as such. A
generic imported input is one whose price has no stochastic volatility: an AR(1)
in the log price with a fixed innovation variance. Matching it to oil on
persistence and unconditional volatility pins it to the constant-volatility
economy of Section 6.2.

\begin{table}[h]
\centering\footnotesize
\begin{tabular}{lccc}
\toprule
economy & persistence of $\log p$ & sd$(\log p)$ & mean $p$ \\
\midrule
full (oil, stochastic volatility)    & 0.94751 & 0.44065 & 1.24294 \\
constant volatility (generic input)  & 0.94772 & 0.44025 & 1.24211 \\
difference                           & $+0.02\%$ & $-0.09\%$ & $-0.07\%$ \\
\bottomrule
\end{tabular}
\end{table}

The comparison the referee asks for is therefore the full economy against the
constant-volatility economy, which the decomposition already reports: of the
200.7 basis points by which oil price risk raises the mean spread, 13.4 is
attributable to the time variation in volatility and the remaining
\textbf{187 basis points --- 93\% --- is what a generic imported input with
oil's persistence and volatility would deliver}. On the default rate the
time-varying part is not distinguishable from zero. Welfare splits the other way:
removing all oil price risk is worth 0.34\% to 0.46\% of consumption and removing
only the time variation in volatility 0.19\% to 0.22\%, so roughly half of the
welfare cost does come from the stochastic volatility.

What is specific to oil is therefore not the structure but which input satisfies
its premise. For the economies in our sample oil is 9 to 10\% of GDP, its log
price has a quarterly persistence of 0.95, and its unconditional standard
deviation is 0.44. Few imported inputs are simultaneously that large and that
volatile.

\subsection*{Is the effect stronger where oil dependence is greater?}

\begin{quote}\itshape
One natural empirical test would be to examine whether the effect of global oil
price volatility is stronger in countries with greater oil import dependence.
\end{quote}

We split the sample into thirds of the country-level oil import share and
estimate a separate volatility slope within each, which imposes no functional
form on how the effect varies with dependence.

\begin{table}[h]
\centering\footnotesize
\begin{tabular}{lccc}
\toprule
third of the oil import share distribution & volatility slope & s.e. & $t$ \\
\midrule
least dependent ($\leq 1.83\%$ of GDP) & $+0.0108$ & 0.0052 & $+2.06$ \\
middle ($\leq 2.95\%$)                 & $+0.0106$ & 0.0046 & $+2.29$ \\
most dependent ($> 2.95\%$)            & $+0.0085$ & 0.0043 & $+1.98$ \\
\bottomrule
\end{tabular}
\end{table}

The spread response to oil price volatility is present in every third of the
distribution and of similar magnitude throughout; the three slopes are not
statistically distinguishable from one another. The effect is therefore not
confined to the most exposed economies, although the predicted gradient is not
visible in this cross-section either. With 32 countries whose import shares span
1.1 to 10.4 per cent of GDP, the sample cannot resolve differences of the size
the model implies, so we present the cross-sectional prediction as untested
rather than confirmed.

\subsection*{Robustness to generic input-cost volatility}

\begin{quote}\itshape
Additional robustness could control for broader imported-input or producer-price
volatility to assess whether the oil volatility coefficient contains information
beyond a generic global input cost shock.
\end{quote}

The comment names two things, imported-input volatility and producer-price
volatility, and we address both. For the first we use the IMF global metals
price index, metals being an imported intermediate input rather than a
consumption good; for the second a producer price index for \emph{intermediate
demand}, which is the producer-price measure defined on inputs rather than on
finished output. We also add the IMF non-fuel commodity index as the broadest
available measure of a generic global input cost. Each is fitted with the same
estimator, design and seed as the paper's oil volatility ---
\texttt{svsample(log(real), designmatrix = "ar1")} --- on quarterly data
deflated by the same CPI, so every series is measured like for like.

\begin{table}[h]
\centering\footnotesize
\setlength{\tabcolsep}{5pt}
\begin{tabular}{lccccc}
\toprule
control added & oil vol. & $t$ & the control & $t$ & corr.\ with oil \\
\midrule
\emph{none: the paper's specification} & $+0.0100$ & $+2.29$ & --- & --- & --- \\
metals (industrial inputs)             & $+0.0102$ & $+2.42$ & $-0.0013$ & $-0.14$ & $+0.33$ \\
PPI, intermediate demand               & $+0.0085$ & $+1.66$ & $+0.0154$ & $+0.51$ & $+0.68$ \\
non-fuel commodities                   & $+0.0077$ & $+1.74$ & $+0.0158$ & $+1.51$ & $+0.49$ \\
\bottomrule
\end{tabular}
\end{table}

\textbf{The oil coefficient survives all three.} Against metals it is
unchanged --- $+0.0102$ against $+0.0100$, and slightly better determined ---
while the metals coefficient is zero. That is the most informative of the three
tests, because the two volatility series are correlated at only $+0.33$, so the
regression can separate them, and the explanatory power stays with oil. Against
the intermediate-demand PPI and the non-fuel index the oil coefficient is
attenuated to a $t$ of 1.7, but the controls are themselves insignificant and the
correlations are higher: that is two collinear series splitting one signal rather
than a control absorbing it. Each pair is jointly significant.

We also ran four further indices. The coefficient is displaced by a food price
index and by a \emph{finished-goods} producer price index --- but neither is an
imported intermediate input: food is a consumption commodity and a finished-goods
index prices output rather than inputs. They are reported for completeness, and
we do not read them as bearing on the question put here.

Two qualifications belong with this result. The first is that the robustness
above is a full-sample property. Splitting at the financial crisis, oil volatility is
insignificant before 2008 against every control and significant after 2009
against every control, with a larger coefficient of $+0.013$ to $+0.018$. The
The second is that the later window is short, about twenty quarters, so we rely
on the consistency of the pattern across controls rather than on individual $t$
statistics.

\subsection*{Suggested manuscript wording}

\begin{quote}
Several papers study commodity prices in sovereign default models. Lopez-Martin
et al.\ (2019) consider stochastic commodity revenues and hedging, and Hamann
et al.\ (2018) and Esquivel (2020) study oil in exporting economies; in all of
these the commodity is a source of revenue. Seoane (2019) studies stochastic
volatility in the endowment and Johri, Khan and Sosa-Padilla (2022) stochastic
volatility in the world interest rate, the latter documenting the state
dependence in volatility and debt that our Figures 3 and 5 also display. What is
new here is the combination of oil importers, an oil price process whose
persistence and stochastic volatility are estimated rather than assumed, and a
transmission running through the cost side of production. The mechanism is in
principle a general imported-input mechanism: a generic input matched to oil on
persistence and unconditional volatility delivers 187 of the 201 basis points by
which oil price risk raises the mean spread. What makes oil the relevant
application is that it is both large enough, at 9 to 10 per cent of GDP, and
volatile enough, at an unconditional standard deviation of 0.44 in logs, for the
mechanism to bind; and the estimated spread response to oil price volatility is
unchanged when volatility in industrial input prices, in intermediate-demand
producer prices, or in non-fuel commodities generally is controlled for.
\end{quote}
